% Generated draft: requires mathematical and source review.
\documentclass[11pt]{article}
\usepackage[margin=25mm]{geometry}
\usepackage{fontspec}
\setmainfont{Latin Modern Roman}
\usepackage{amsmath,amssymb,mathtools,unicode-math}
\setmathfont{Latin Modern Math}
\usepackage{enumitem,needspace,xurl,hyperref}
\setlength{\parindent}{0pt}
\setlength{\parskip}{5pt}
\setlength{\emergencystretch}{3em}
\newcommand{\N}{\mathbb{N}}
\newcommand{\Z}{\mathbb{Z}}
\newcommand{\Q}{\mathbb{Q}}
\newcommand{\R}{\mathbb{R}}
\newcommand{\C}{\mathbb{C}}
\newcommand{\F}{\mathbb{F}}
\newcommand{\A}{\mathbb{A}}
\newcommand{\PP}{\mathbb{P}}
\newcommand{\E}{\mathbb{E}}
\newcommand{\eps}{\varepsilon}
\newcommand{\dd}{\,\mathrm{d}}
\begin{document}
\section*{UnsolvedMath ID: 30002388}
\textbf{Source problem code:} catalogue entry
\subsection{Definitions and mathematical statement}
Write \(m_{\mathrm{GBC}}(M,g)\) for the Gauss--Bonnet--Chern mass referred to in the supplied record. This notation is only a placeholder: the record gives neither its defining formula nor its normalization, and does not specify a curvature order or dimension range. It also does not define the intended asymptotic flatness, decay, regularity, or integrability conditions.

\textbf{Source-level question.} Do the Positive Mass Theorem and the Penrose Inequality for the Gauss--Bonnet--Chern mass hold for asymptotically flat manifolds?

The nonnegativity part of the intended positive mass statement would have the schematic form
\[
 m_{\mathrm{GBC}}(M,g)\geq 0
\]
for every manifold satisfying the intended hypotheses. Those hypotheses are not supplied, so this is not a conjecture asserting nonnegativity for every asymptotically flat manifold without further conditions. The record does not specify completeness, permissible boundary or ends, a curvature sign condition, or an equality and rigidity conclusion.

The intended Penrose inequality asks for a lower bound on this mass in terms of suitable horizon or boundary geometry. The record supplies no formula for that bound, dimensional constant, exponent, horizon or boundary assumptions, or equality case. Consequently, a precise quantified Penrose inequality cannot be recovered from this record alone. Both questions are retained at their supplied level of generality and require the missing definitions and hypotheses before they can be stated as definite mathematical conjectures.
\subsection{Short English statement}
The problem asks whether positive mass and Penrose inequalities extend to the Gauss--Bonnet--Chern mass on general asymptotically flat manifolds. The supplied record leaves the mass definition, necessary geometric hypotheses, and precise Penrose bound unspecified, so this draft preserves the question and explicitly records those gaps.
\subsection{Sources}
\begin{itemize}
\item Local UnsolvedMath registry, permanent ID \(30002388\), source problem code \texttt{catalogue entry}, title: \emph{Gauss--Bonnet--Chern Positive Mass and Penrose Inequalities}. The record belongs to the Oberwolfach Reports Open Problems collection and records proposed year 2013; no proposer is named.
\item Supplied source citation: \emph{Partial Differential Equations (2014)}, Oberwolfach Reports, DOI 10.4171/owr/2013/40, pp.\ 2259--2319. This is also the supplied literature source labeled ``original OWR report'': \url{https://doi.org/10.4171/owr/2013/40}. The record describes Oberwolfach Reports as documenting mathematical workshops and their open-problem discussions.
\item Supplied literature source labeled ``Ge--Wang--Wu primary paper'': \url{https://arxiv.org/abs/1211.3645}. No theorem number or precise supporting formulation is supplied.
\item Literature assessment dated exactly \texttt{2026-08-21}: the supplied curation labels the problem open and states that positive mass and Penrose results are known for graphs and conformally flat manifolds, while the conjectured general asymptotically flat theorem remains unresolved. This is supplied dated curation, not independently verified here; its time-sensitive status and cited claims require independent checking.
\item The supplied record labels the statement \texttt{corrected\_verified} and reports that its mathematical content was checked against the official Oberwolfach report, with notation, context, and question form normalized without changing the source-level claim. This is a recorded source-verification claim, not verification performed for this draft.
\end{itemize}
\end{document}
